% Generated draft: requires mathematical and source review.
\documentclass[11pt]{article}
\usepackage[margin=25mm]{geometry}
\usepackage{fontspec}
\setmainfont{Latin Modern Roman}
\usepackage{amsmath,amssymb,mathtools,unicode-math}
\setmathfont{Latin Modern Math}
\usepackage{enumitem,needspace,xurl,hyperref}
\setlength{\parindent}{0pt}
\setlength{\parskip}{5pt}
\setlength{\emergencystretch}{3em}
\newcommand{\N}{\mathbb{N}}
\newcommand{\Z}{\mathbb{Z}}
\newcommand{\Q}{\mathbb{Q}}
\newcommand{\R}{\mathbb{R}}
\newcommand{\C}{\mathbb{C}}
\newcommand{\F}{\mathbb{F}}
\newcommand{\A}{\mathbb{A}}
\newcommand{\PP}{\mathbb{P}}
\newcommand{\E}{\mathbb{E}}
\newcommand{\eps}{\varepsilon}
\newcommand{\dd}{\,\mathrm{d}}
\begin{document}
\section*{UnsolvedMath ID: 6700082}
\textbf{Source problem code:} catalogue entry
\subsection{Definitions and mathematical statement}
Let $(X,d)$ be a compact $n$-dimensional pseudomanifold with the source's Hilbert-volume structure, and let $S^n$ carry its unit round metric and Riemannian volume. The proposed small-ball hypothesis is
\[
\operatorname{Vol}_{\mathrm H}(B_X(x,r))
<\operatorname{Vol}_{S^n}(B_{S^n}(q,r))
\qquad(x\in X,\ 0<r\leq\varepsilon_0),
\]
where the right side is independent of $q\in S^n$. The proposed conclusion is that every $\lambda$-Lipschitz map $f:X\to S^n$ is null-homotopic. This means there exists a continuous homotopy from $f$ to a constant map; it is stronger than merely asserting degree zero and does not assert that $X$ itself is contractible.

The source suggests, tentatively, starting values $\varepsilon_0=\pi/4$ and $\lambda\leq1/2$. A $\lambda$-Lipschitz map satisfies $d_{S^n}(f(x),f(y))\leq\lambda d(x,y)$ for all $x,y\in X$. The stacked-fraction values are rendered here in their natural $\pi/4$ and $1/2$ interpretation, which requires checking against the normalized extraction. No universal conclusion for unspecified alternative constants is inferred.

Hilbert volume and the admissible metric/pseudomanifold category are not defined in the record. Hilbert volume cannot simply be replaced by Hausdorff measure or an arbitrary measure, since the ball hypothesis depends on its normalization and relation to the metric. The dimension and boundary/singularity conventions are also absent.

The author's qualification is essential: the source calls conjectures [79]--[84] ``on the side of wishful thinking'' and says that something opposite may be true. It ends by introducing a following example that is not included. This is a deliberately tentative synthetic comparison proposal, not an established non-Riemannian extension of Llarull rigidity.
\subsection{Short English statement}
If every sufficiently small Hilbert-volume ball is smaller than its spherical counterpart, must all sufficiently short maps to the sphere be null-homotopic? The source tentatively suggests radii up to $\pi/4$ and Lipschitz bound $1/2$, while explicitly doubting the proposal.
\subsection{Sources}
\begin{itemize}
\item Local UnsolvedMath registry, permanent ID 6700082, exact source code catalogue entry, ``Non-Riemannian $\varepsilon$-Llarull''; proposed by Misha Gromov, 2017. Registry status open, classification OPEN-TRIAGE. Source printed tag [?84], occurrence 85.
\item Misha Gromov, \emph{101 Questions, Problems and Conjectures around Scalar Curvature} (2017), PDF page 83, \url{https://www.ihes.fr/~gromov/wp-content/uploads/2018/08/101-problemsOct1-2017.pdf}. Supplied normalized pdf-text extraction accessed July 29, 2026; incomplete author-hosted PDF. No PDF was accessed here.
\item The supplied summary describes a synthetic Hilbert-volume analogue of Llarull's rigidity and reports no resolution located. No precise later reference is supplied. The source labels [79]--[84] are its conjecture identifiers, not bibliography entries supplied here.
\item Provenance flags current status and statement-reuse terms NEEDS\_REVIEW. The author's explicit wishful-thinking and possible-opposite warnings are retained.
\end{itemize}
\end{document}
